\documentclass[11pt,a4paper]{article}
\usepackage[utf8]{inputenc}
\usepackage{amsmath,amssymb,amsfonts}
\usepackage{geometry}
\geometry{margin=1in}

\title{\textbf{Formulación de una Acción Unificada EMD y su Reducción al Límite Clásico Electrostático}}
\author{Tu Nombre / Autor}
\date{\today}

\begin{document}

\maketitle

\begin{abstract}
Este documento presenta una propuesta formal de Acción unificada ($\mathcal{S}$) que combina la Relatividad General de Einstein, el electromagnetismo de Maxwell y la teoría de campos fermiónicos de Dirac. Se analiza la consistencia dimensional del modelo y se demuestra su reducción algebraica al límite clásico de la Ley de Coulomb bajo condiciones estáticas.
\end{abstract}

\section{Introducción y Marco Teórico}
La unificación de las interacciones fundamentales requiere un formalismo basado en la teoría de campos y la geometría diferencial. Proponemos la siguiente densidad de Acción acoplada en un espacio-tiempo curvo:

\begin{equation}
\mathcal{S} = \int d^4x \sqrt{-g} \left[ \frac{c^4}{16\pi G} R - \frac{1}{4} F_{\alpha\beta} F^{\alpha\beta} + \bar{\psi}(i \hbar \gamma^\mu D_\mu - m c)\psi \right]
\end{equation}

Donde $\sqrt{-g}$ representa el determinante de la métrica espaciotemporal, $R$ el escalar de curvatura de Ricci, $F_{\alpha\beta}$ el tensor electromagnético y $\psi$ el espinor de Dirac con su respectiva derivada covariante de gauge $D_\mu$.

\section{Reducción al Límite Clásico (Ley de Coulomb)}
Al considerar un sistema en el régimen estático local ($\vec{v} = 0$, $\gamma = 1$, curvatura gravitacional despreciable $R \to 0$) y en métrica plana de Minkowski, el término de interacción electromagnética se desacopla de la métrica general, recuperando la ley de fuerzas electrostáticas de Coulomb:

\begin{equation}
F_e = \frac{1}{4\pi \varepsilon_0} \frac{|q_1 q_2|}{r^2}
\end{equation}

\section{Evidencia y Validación Numérica}
Se evaluó el modelo utilizando un caso de prueba controlado con cargas puntuales:
\begin{itemize}
    \item Carga 1 ($q_1$): $-2 \, \mu\text{C}$
    \item Carga 2 ($q_2$): $+4 \, \mu\text{C}$
    \item Distancia ($r$): $0.04 \, \text{m}$
\end{itemize}

Mediante el cálculo estándar, se obtiene una fuerza de $F_e = 44.9377 \, \text{N}$ y un trabajo/energía asociada de $E = 1.7975 \, \text{J}$, validando la consistencia del marco clásico de la teoría.

\section{Conclusión}
La estructura de la Acción EMD demuestra ser matemáticamente coherente en sus dimensiones de densidad de energía y permite recuperar de forma natural los límites clásicos conocidos al aplicar restricciones locales.

\end{document}
